\documentclass[12pt]{article}

% Layout.
\usepackage[top=1in, bottom=0.75in, left=1in, right=1in, headheight=1in, headsep=6pt]{geometry}

% Fonts.
\usepackage{mathptmx}
\usepackage[scaled=0.86]{helvet}
\renewcommand{\emph}[1]{\textsf{\textbf{#1}}}

% TiKZ.
\usepackage{tikz, pgfplots}

\usetikzlibrary{calc,trees,positioning,arrows,fit,shapes,through, backgrounds}
\usetikzlibrary{patterns}

\usetikzlibrary{decorations.markings}
\usetikzlibrary{arrows}

\usepackage{pgfplots}


% Misc packages.
\usepackage{amsmath,amssymb,latexsym}
\usepackage{graphicx}
\usepackage{array}
\usepackage{xcolor}
\usepackage{multicol}
\usepackage{fancyhdr}
\pagestyle{fancy}


% Misc.
\renewcommand{\d}{\displaystyle}
\newcommand{\ds}{\displaystyle}
\newcommand{\ul}[1]{\underline{#1}}
\def\bc{\begin{center}}
\def\ec{\end{center}}
\def\be{\begin{enumerate}}
\def\ee{\end{enumerate}}

\newcommand{\ans}[1][1in]{\rule{#1}{.5pt}}


\lhead{Math F113X: Quiz 4}
\rhead{October 3, 2026}

\begin{document}

\vspace{1cm}
\strut

\noindent \textbf{Name:} \hrulefill \quad  score:\ans[1cm] \ / 10 \\


\noindent There are 10 points possible on this quiz. No aids (book, notes, etc.)
are permitted. You may use a calculator.  {\bf Show all work and supporting calculations for full credit. Explain how you get your answers.}



\be

\item (6 points) Mary and Robert are going to divide their shared kitchen gadgets, an instant pot and a microwave, using the method of sealed bids. \emph{ Fill in the blanks below. }
\begingroup

%\setlength{\tabcolsep}{10pt} % Default value: 6pt
\renewcommand{\arraystretch}{1.5} % Default value: 1
\begin{tabular}{|p{3in} || >{\centering\arraybackslash}m{1.5in} |>{\centering\arraybackslash}m{1.5in} |}
\hline
& Mary & Robert \\ \hline \hline
instant pot & \$100 & \$90 \\ \hline
microwave & \$50 & \$70 \\ \hline \hline
total bid & \$150 & \$160 \\ \hline
 What is a \emph{fair share} for each person?& & \\ \hline
 \emph{award} (who gets what item)& &  \\ \hline
{award value}&  &   \\ \hline
  How much money does each person pay in or receive from the holding pile? &&  \\
 Make sure to indicate whether the person \emph{pays} or \emph{receives}.&&\\  \hline
 Determine the surplus. &\multicolumn{2}{c|}{\quad} \\[12 pt] \hline
 Determine each person's fair share of the  surplus. &&\\[12 pt] \hline
\end{tabular}

\vfill
 Determine the Final Allocation by filling out each blank below.\\


\renewcommand{\arraystretch}{2} 
\begin{tabular}{ | l || >{\centering\arraybackslash}p{4cm}| >{\centering\arraybackslash}p{4cm} |>{\centering\arraybackslash}p{4cm} |}
\hline
\multicolumn{4}{| c |}{Summary of Final Allocation}\\
\hline
person& items & cash & paid or received? \\

\hline

Mary&&&\\

\hline

Robert &&&\\
\hline
\end{tabular}
\ee


\newpage

\item (2 points) Describe a situation that can be modeled with a  graph.% What do the vertices represent? What do the edges represent? What do the weights represent?

What do the \emph{vertices} represent?

\vfill

What do the \emph{edges} represent? 

\vfill



\item (2 points) Consider the following graph:

\makeatletter
\newcommand{\Letter}[1]{\@Alph{#1}}
\makeatother

\begin{center}
\begin{tikzpicture}[baseline=(current bounding box.center),vertex/.style = {circle, draw, inner sep=2pt}]

\foreach \i in {1,...,5}{\node[vertex] (\Letter{\i}) at (360*\i/5:2) {\Letter{\i}};}
%\node[vertex, right = of E] (F) {F};
\node[vertex, right = 3 cm of E] (G) {S};
\foreach \i/\j in {A/D,  B/C, C/E, E/B}{\draw (\i) -- (\j);}
\draw (E) edge[bend left = 30] (G);
\draw (E) edge[bend right = 30] (G);
\draw (G) to[out = 45, in = -45, looseness = 3, distance = 2cm] (G);
\end{tikzpicture}
\end{center}

\be
\item How many vertices does this graph have? \ans
\item How many edges does this graph have? \ans
\item What is the degree of vertex $S$? \ans
\item Explain why this graph is not connected.

\vspace{1in}
\ee


	

	

 \ee
\end{document}